\chapter{Perfectly Secret Encryption}

\noindent
\textbf{2.1} \hspace{1em} Prove that, by redefining the key space, we may assume that the key-generation algorithm \textbf{Gen} chooses a uniform key from the key space, without changing Pr[C= c|M= m] for any m,c.
\vspace{1em}

Motivation: It's is easier to think matematically with a Gen that generates keys at an uniform distribution.

So, let's use a new key $\omega$ that is uniformly distributed over a new key space $\Omega$. And replace the \textbf{Gen() = k}  by the uniform \textbf{GenU() = w}. We also define \textbf{Map(w) = k} in a way that, the distribuiton of \textit{K}  = Map(Gen(W)) is the same as the original.

As an example, let's say \textit{K} = \{0, 1\} with $\Pr[K = 0] = 0.75$ and $\Pr[K = 1] =  0.25$. We than choose \textbf{W} = \{00, 01, 10, 11\} with equal probability (\textit{uniform distribution}), and make Map(w) = 1 if w = 11, otherwise Map(w) = 0.

With the new keyspace we have:

\[
\Pr[C = c \mid M = m] = \sum_{w \in \mathcal{W}} \Pr[\text{GenU()} = w] \cdot \Pr[\text{Enc}_{map(w)}(m) = c]
\]

Also notice that for each \( k \) in the original scheme can be produced by one or more \( \omega \) in the new scheme. That is:
\[
\Pr[\text{Gen()} = k] = \sum_{\omega : \text{Map}(\omega) = k} \Pr[GenU() = \omega] = \sum_{\omega} \Pr[GenU() = \omega] \cdot \Pr[\text{Enc}_{\text{Map}(\omega)}(m) = c]
\]


Let's prove that from the original key scape.

\[
\Pr[C = c \mid M = m] = \sum_{k \in \mathcal{K}} \Pr[\text{Gen()} = k] \cdot \Pr[\text{Enc}_k(m) = c]
\]

\[
 = \sum_{k \in \mathcal{K}} \left( \sum_{w : \text{map}(w) = k} \Pr[GenU() = w] \right) \cdot \Pr[\text{Enc}_k(m) = c]
\]

Flatting the summation we get the new schema:

\[
= \sum_{w \in \mathcal{W}} \Pr[GenU() = w] \cdot \Pr[\text{Enc}_{\text{map}(w)}(m) = c]
\]

\vspace{1em}
\noindent
\textbf{2.2} \hspace{1em} Prove that, by redefining the key space as well as the encryption algorithm, we may assume that encryptation is deterministic without changing Pr[C= c|M= m] for any m,c.
\vspace{1em}

Motivation: It's is easier to think matematically with a determistic Enc instead of randomized.

If $Enc_k(m)$ is randomized, it generates $c \in C_{m, k} = \{c_0, c_1, ..., c_l\}$ chiphered messages. If we compute then randomized part $\omega$, separately, we can rewrite Enc as  $Enc_k(m, \omega)$.

We can now move out the internal random part and make it as part of the Key. So, the Key is now $k\omega$. Then:

Gen() = k becomes $Gen'() = k\omega$

$Enc_k(m)$ becomes $Enc_{k\omega}'(m) := Enc_k(m, \omega)$

$Dec_k(c)$ becomes  $Dec_{k\omega}'(c)$ that is equal to the original $Dec_k(c)$, because Dec is deteministic and doesn't depend on $\omega$.

From original schema:

\[
\Pr[C = c \mid M = m] = \sum_{k} \Pr[k]  \sum_{\omega} \Pr[\omega] \cdot \Pr[\text{Enc}_k(m, \omega) = c]
\]

In new schema:

\[
\Pr[C = c \mid M = m] = \sum_{k\omega} \Pr[k\omega] \cdot \Pr[\text{Enc}'_{k\omega}(m) = c]
\]

\[
\Pr[C = c \mid M = m] = \sum_{k\omega} \Pr[k\omega] \cdot \Pr[\text{Enc}_k(m, \omega) = c]
\]

Now, since  $\Pr[k\omega]=\Pr[k]\cdot\Pr[\omega]$ (\textit{i.e.} they are independent), the expression becomes:

\[
\Pr[C = c \mid M = m] = \sum_{k} \sum_{\omega} \Pr[k]  \Pr[\omega] \cdot \Pr[\text{Enc}_k(m\omega) = c]
\]

\[
\Pr[C = c \mid M = m] = \sum_{k} \Pr[k] \sum_{\omega} \Pr[\omega] \cdot \Pr[\text{Enc}_k(m\omega) = c]
\]

Which is the same as the original.

\vspace{1em}
\noindent
\textbf{2.3} \hspace{1em} Prove or refute: An encryption scheme with message space $\mathcal{M}$ is perfectly secret if and only if for every probability distribution on $\mathcal{M}$ and every $c0,c1 \in \mathcal{C}$ we have $\Pr[C= c0] = \Pr[C= c1]$.
\vspace{1em}

This is false as we can see from encryptation scheme bellow.

\begin{align*}
	\mathcal{M} &= m_0 = \{0,1\} \\
	\mathcal{K} &= k_0k_1r \in \{0,1\}^3 = \{000,\,001,\,010,...,\,110,\,111\} \\
	\mathcal{C} &=  c_0c_1 \in \{0,1\}^2 = \{00,\,01,\,01,\,111\} \\
	Enc_k(m) &=  c_0c_1 \text{ where } c_0 = m \oplus k_0 \oplus r \text{ and } c_1 = r \oplus k_1 \\
    Dec_k(m) &= m_0 = c_0 \oplus k_0 \oplus r \text{ where } r = c_1 \oplus k_1 \\
\end{align*}

Then, let's assume the $Pr[M = 0] = 0.5, Pr[k_0 = 0] = 0.5, Pr[k_1 = 0.25] = 0.5, Pr[r = 0] = 0.25$.

\begin{tabular}{|c|c|c|c|c|c|c|c|}
	\hline
	$m$ & $k_0$ & $k_1$ & $r$ & $m \oplus k_0 \oplus r$ & $k_1 \oplus r$ & $C$ & $\Pr(C)$ \\
	\hline
	0 & 0 & 0 & 0 & 0 & 0 & 00 & 0.015625 \\
	0 & 0 & 0 & 1 & 1 & 1 & 11 & 0.046875 \\
	0 & 0 & 1 & 0 & 0 & 1 & 01 & 0.046875 \\
	0 & 0 & 1 & 1 & 1 & 0 & 10 & 0.140625 \\
	0 & 1 & 0 & 0 & 1 & 0 & 10 & 0.015625 \\
	0 & 1 & 0 & 1 & 0 & 1 & 01 & 0.046875 \\
	0 & 1 & 1 & 0 & 1 & 1 & 11 & 0.046875 \\
	0 & 1 & 1 & 1 & 0 & 0 & 00 & 0.140625 \\
	1 & 0 & 0 & 0 & 1 & 0 & 10 & 0.015625 \\
	1 & 0 & 0 & 1 & 0 & 1 & 01 & 0.046875 \\
	1 & 0 & 1 & 0 & 1 & 1 & 11 & 0.046875 \\
	1 & 0 & 1 & 1 & 0 & 0 & 00 & 0.140625 \\
	1 & 1 & 0 & 0 & 0 & 0 & 00 & 0.015625 \\
	1 & 1 & 0 & 1 & 1 & 1 & 11 & 0.046875 \\
	1 & 1 & 1 & 0 & 0 & 1 & 01 & 0.046875 \\
	1 & 1 & 1 & 1 & 1 & 0 & 10 & 0.140625 \\
	\hline
	\multicolumn{7}{|r|}{sum} & 1 \\
	\hline
\end{tabular}

\vspace{1em}
Notice it is still perfectly secret

\begin{tabular}{|c|c||c|c|}
	\hline
	\multicolumn{2}{|c||}{Pr$[C \mid M=0]$} & \multicolumn{2}{c|}{Pr$[C \mid M=1]$} \\
	\hline
	$C$ & Prob & $C$ & Prob \\
	\hline
	00 & 0.15625 & 00 & 0.15625 \\
	01 & 0.09375 & 01 & 0.09375 \\
	10 & 0.15625 & 10 & 0.15625 \\
	11 & 0.09375 & 11 & 0.09375 \\
	\hline
	sum & 0.5 & sum & 0.5 \\
	\hline
	\multicolumn{3}{|r|}{sum of sums} & 1 \\
	\hline
\end{tabular}

\vspace{1em}
But not all $c \in C$ has same probability.

\begin{tabular}{|c|c|}
	\hline
	\multicolumn{2}{|c|}{Pr$[C = c_0] \neq$ Pr$[C = c_1]$} \\
	\hline
	$C$ & Prob \\
	\hline
	00 & 0.3125 \\
	01 & 0.1875 \\
	10 & 0.3125 \\
	11 & 0.1875 \\
	\hline
	sum & 1 \\
	\hline
\end{tabular}

\newpage
\noindent
\textbf{2.4} \hspace{1em} Prove or refute: For every perfectly secret encryption scheme it holds that for every distribution on the message space $\mathcal{M}$, every $m, m^\prime, \in \mathcal{M}$ and every $c \in \mathcal{C}$:

$\Pr[\mathcal{M}= m \mid \mathcal{C}= c] = \Pr[\mathcal{M}= m^\prime \mid \mathcal{C}= c]$
\vspace{1em}

From $\Pr[\mathcal{M}= m \mid \mathcal{C}= c] = \Pr[\mathcal{M}= m]$ (Lemma 2.3 - definition of perfect secret)
\[
\frac{\Pr[\mathcal{M}= m \mid \mathcal{C}= c]}{\Pr[\mathcal{M}= m]} = 1
\]
Bayes’ Theorem gives:
\[
\Pr[M = m \mid C = c] = \frac{\Pr[C = c \mid M = m] \cdot \Pr[M = m]}{\Pr[C = c]}
\]
Therefore:
\[
\frac{\Pr[M = m \mid C = c]}{\Pr[M = m]} = \frac{\Pr[C = c \mid M = m]}{\Pr[C = c]} = 1
\]
and:
\[
\frac{\Pr[M = m^\prime \mid C = c]}{\Pr[M = m^\prime]} = \frac{\Pr[C = c \mid M = m^\prime]}{\Pr[C = c]} = 1
\]
then:
\[
\Pr[C = c \mid M = m] = \Pr[C = c \mid M = m^\prime]
\]

\newpage
\noindent
\textbf{2.5} \hspace{1em}Prove that in Definition 2.6 we may assume $\mathcal{A}$ is deterministic without loss of generality.
\vspace{1em}

From the experiment. The encryptation scheme $\Pi = (\text{Gen}, \text{Enc}, \text{Dec})$, and the randoness of $k \leftarrow \text{Gen}$ and the choice of what message will be encrypted, \textit{i.e}, the choice of bit $b$, is not under the control of adversary $\mathcal{A}$.

The adversary $\mathcal{A}$, can only choose two messages $m_0, m_1 \in \mathcal{M} \mid m_0 \ne m_1$, and the output bit $b^\prime$.

So, let's fix the experiment, \textit{i.e.} the part adversary $\mathcal{A}$ has not control of.

Then, consider $\mathcal{A}$ is random. Since there are a finite set of combinaitions of  $m_0, m_1 \in \mathcal{M} \mid m_0 \ne m_1$ and $b^\prime$, there is one of these.

For a random $\mathcal{A}$, we have:

\[
\Pr[\text{Priv}^{\text{eav}}_{A, \Pi} = 1] = \frac{1}{2}
\]

Now, let's take the random part out of $\mathcal{A}$, with  $r \leftarrow \mathcal{R}$, and rewrite a deterministic one as  $\mathcal{A}^\prime(r)$. Now we have:

\[
\Pr[\text{Priv}^{\text{eav}}_{A, \Pi} = 1] = \sum_{r \in \mathcal{R}} \Pr[\text{Priv}^{\text{eav}}_{A^\prime(r), \Pi} = 1] \cdot \Pr[R = r] = \mathbb{E}_r[\Pr[\text{Priv}^{\text{eav}}_{A^\prime(r), \Pi} = 1]]
\]

\textit{I.e.} it is the average or expected value of $\Pr[\text{Priv}^{\text{eav}}_{A^\prime(r), \Pi} = 1]$ with $r \in R$.

So, there is a $r^\prime$ where:

$\Pr[\text{Priv}^{\text{eav}}_{A^\prime(r^\prime), \Pi} = 1] \ge \mathbb{E}_r[\Pr[\text{Priv}^{\text{eav}}_{A^\prime(r), \Pi} = 1]] = \Pr[\text{Priv}^{\text{eav}}_{A, \Pi} = 1] = \frac{1}{2}$


But, we can't alllow $\Pr[\text{Priv}^{\text{eav}}_{A^\prime(r^\prime), \Pi} = 1] > \frac{1}{2}$ for any $r$, because, if we do so, then there would be some randoness that would choose just those good $r$s and then $\Pr[\text{Priv}^{\text{eav}}_{A, \Pi} = 1] > \frac{1}{2}$. And we can't allow any r so that  $\Pr[\text{Priv}^{\text{eav}}_{A^\prime(r^\prime), \Pi} = 1] < \frac{1}{2}$, because it would make $ \mathbb{E}_r[\Pr[\text{Priv}^{\text{eav}}_{A^\prime(r), \Pi} = 1]] < \frac{1}{2}$.

Then, finnaly:

$\Pr[\text{Priv}^{\text{eav}}_{A^\prime(r^\prime), \Pi} = 1] = \Pr[\text{Priv}^{\text{eav}}_{A, \Pi} = 1] = \frac{1}{2}$


\newpage
\noindent
\textbf{2.6} \hspace{1em} Prove Lemma 2.7.
\vspace{1em}

First le't prove that if it is perfect secret, then $\Pr[\text{Priv}^{\text{eav}}_{A, \Pi} = 1]$ is $\frac{1}{2}$ .

\vspace{1em}

$\Pr[\text{Priv}^{\text{eav}}_{A, \Pi}\ = 1] = \Pr[b^\prime = b]$

$= \Pr[b^\prime = 0 \mid M = m_0]\Pr[M = m_0] + \Pr[b^\prime = 1 \mid M = m_1]\Pr[M = m_1]$

\vspace{1em}

$\mathcal{A}$ is determistic, then, let's assume that if $c \in C_0$  it outputs $b^\prime = 0$ if $c \in C_1$  it outputs $b^\prime = 1$.

$= \Pr[c \in C_0 \mid M = m_0]\Pr[b = 0] + \Pr[c \in C_1 \mid M = m_1]\Pr[b = 1]$

$= \Pr[c \in C_0 \mid M = m_0]\cdot\frac{1}{2} + \Pr[c \in C_1 \mid M = m_1]\cdot\frac{1}{2}$

$=  \sum_{c \in C_0} \Pr[C = c \mid M = m_0]\cdot\frac{1}{2} +  \sum_{c \in C_1} \Pr[C = c \mid M = m_1]\cdot\frac{1}{2}$

\vspace{1em}

If it is perfectly secret, then (from Definition 2.3 and Bayes' theorem):  $\Pr[C = c \mid M = m] = \Pr[C = c]$. Therefore:

\vspace{1em}

$=  \sum_{c \in C_0} \Pr[C = c]\cdot\frac{1}{2} +  \sum_{c \in C_1} \Pr[C = c]\cdot\frac{1}{2}$

$= \frac{1}{2} ( \sum_{c \in C_0} \Pr[C = c] +  \sum_{c \in C_1} \Pr[C = c] )$

\vspace{1em}

Also notice that $C_0 \cup C_1 = C$ and $C_0 \cap C_1 = \emptyset$. Then:

\vspace{1em}

$= \frac{1}{2}  \cdot \sum_{c \in C} \Pr[C = c] = \frac{1}{2}$

\vspace{1em}

To prove the opposite direction, \textit{i.e}, if $\Pr[\text{Priv}^{\text{eav}}_{A, \Pi} = 1] = \frac{1}{2}$, then it is perfectly secret. We will do contraposite, \textit{i.e}, if \textbf{not} perfectly secret than it is \textbf{not} perfectly indistinguishable.

If \textbf{not} perfectly secret, then, from defintion, there must a $m^\prime$ and $c^\prime$ where:

$\Pr[M = m^\prime \mid C = c^\prime] \ne \Pr[M = m^\prime]$

Also, notice that:

$\sum_{m \in \mathcal{M}} \Pr[M = m] = 1$

$ = \sum_{m \in \mathcal{M} \setminus \{m^\prime\}} \Pr[M = m] + \Pr[M = m^\prime]$ 

And

$\sum_{m \in \mathcal{M}} \Pr[M = m \mid C = c^\prime] = 1$


$= \sum_{m \in \mathcal{M} \setminus \{m^\prime\}} \Pr[M = m \mid C = c^\prime] + \Pr[M = m^\prime \mid C = c^\prime]$

Notice that we \textbf{can't} assume that  for all $m \in \mathcal{M} \setminus \{m^\prime\} \Rightarrow \Pr[M = m \mid C = c^\prime] \ne \Pr[M = m]$

Because, if we do so, in equations above, it would force that $\Pr[M = m^\prime \mid C = c^\prime] = \Pr[M = m^\prime]$, and it is the oposite of our assumption.

So, there must be at least another $m^{\prime\prime}$ that $\Pr[M = m^\prime \mid C = c^\prime] \ne \Pr[M = m^{\prime\prime}]$

\vspace{1em}

Now let's define the deterministica adversary $\mathcal{A}$ (See exercise 2.5), so that $m_0 = m^\prime, m_1 = m^{\prime\prime}$

Then, now let's consider our message space as $\mathcal{M} = \{m_0, m_1\}$.

$\Pr[M = m \mid C = c^\prime] \ne \Pr[M = m] = \frac{1}{2}$

\vspace{1em}

$\Pr[M = m \mid C = c^\prime] = \frac{\Pr[C = c^\prime \mid M = m]\cdot\Pr[M = m]}{\Pr[C = c^\prime]}$  (Bayes' Theorem)

$\Pr[C = c^\prime \mid M = m] = \frac{\Pr[M = m \mid C = c^\prime]}{\Pr[M = m]}\cdot\Pr[C = c^\prime]$

$\Rightarrow \Pr[C = c^\prime \mid M = m] \ne \Pr[C = c^\prime]$

\vspace{1em}

The adversary knows $\mathcal{A}$ that $\Pr[M = m_0 \mid C = c^\prime] \ne \Pr[M = m_1 \mid C = c^\prime]$, and then, one is greater than $\frac{1}{2}$, and the other is smaller. Let's supose $\Pr[M = m_0 \mid C = c^\prime] > \frac{1}{2}  > \Pr[M = m_1 \mid C = c^\prime]$.

Let's also define adversary $\mathcal{A}$ to output $b^\prime = 0$ if $c = c^\prime$, otherwise it outputs $b^\prime \leftarrow \{0, 1\}$. So, that:

$\Pr[\text{Priv}^{\text{eav}}_{A, \Pi}\ = 1] = \Pr[b^\prime = b]$

$= \Pr[b^\prime = 0 \mid M = m_0]\Pr[M = m_0] + \Pr[b^\prime = 1 \mid M = m_1]\Pr[M = m_1]$


$= \frac{1}{2} \cdot ( \Pr[b^\prime = 0 \mid M = m_0] + \Pr[b^\prime = 1 \mid M = m_1])$

let's break down $ \Pr[b^\prime = 0 \mid M = m_0]$

$= \Pr[C = c^\prime \mid M = m_0] + \frac{1}{2} \Pr[C \ne c^\prime \mid M = m_0] $ 

$= \Pr[C = c^\prime \mid M = m_0] + \frac{1}{2} ( 1 -  \Pr[C = c^\prime \mid M = m_0] )$ 

$= \frac{1}{2} + \frac{1}{2} \cdot \Pr[C = c^\prime \mid M = m_0]$ 

let's break down $ \Pr[b^\prime = 1 \mid M = m_1]$

$= \frac{1}{2} \Pr[C \ne c^\prime \mid M = m_1]$ 

So:

$\Pr[\text{Priv}^{\text{eav}}_{A, \Pi}\ = 1] = \frac{1}{2} (\frac{1}{2} + \frac{1}{2} \cdot \Pr[C = c^\prime \mid M = m_0] +  \frac{1}{2} \Pr[C \ne c^\prime \mid M = m_1])$

$= \frac{1}{4} (1 + \Pr[C = c^\prime \mid M = m_0] +  \Pr[C \ne c^\prime \mid M = m_1])$

$\ne \frac{1}{4} (1 + \Pr[C = c^\prime \mid M = m_1] +  \Pr[C \ne c^\prime \mid M = m_1]) = \frac{1}{4} (1 + 1) = \frac{1}{2}$

Therefore:

$\Pr[\text{Priv}^{\text{eav}}_{A, \Pi}\ = 1] \ne \frac{1}{2}$





